% =============================================================================================
%  Préparation de la discussion avec M. Le Maître : sa réponse du 30/10/2026 au mail d'Adrian et
%  sa note « Optimisation de forme d'un foil — VPP et apprentissage adaptatif d'un modèle
%  hydrodynamique gaussien », version 01 du 29/10/2026 (24 p.). Plan de réponse et argumentaire.
%  Note interne (Adrian, Antoine), à ne pas commiter. À lire avec structure_vpp_acceleration.tex.
%  Compilation : tectonic -X compile preparation_reponse_le_maitre.tex --outdir build
% =============================================================================================
\documentclass[11pt,a4paper]{article}
\usepackage{fontspec}
\usepackage[french]{babel}
\usepackage{csquotes}
\usepackage[a4paper,margin=2.0cm,top=1.6cm,bottom=1.8cm]{geometry}
\usepackage{microtype}
\usepackage{amsmath,amssymb,mathtools,bm}
\usepackage{xcolor}
\usepackage{booktabs,tabularx,array}
\usepackage{enumitem}
\usepackage[most]{tcolorbox}
\usepackage{titlesec}
\usepackage[hidelinks,unicode]{hyperref}

\hypersetup{colorlinks=true,linkcolor=blue!50!black,citecolor=green!40!black,
            urlcolor=blue!60!black,pdftitle={Préparer la réponse à M. Le Maître}}
\setlength{\parindent}{0pt}
\setlength{\parskip}{0.4em}
\setlength{\emergencystretch}{3em}
\setlist{itemsep=0.1em,topsep=0.2em,leftmargin=1.3em}
\renewcommand{\arraystretch}{1.12}
\titleformat{\section}{\normalfont\large\bfseries\color{blue!40!black}}{\thesection}{0.6em}{}
\titlespacing*{\section}{0pt}{0.8em}{0.25em}

\tcbset{enhanced,breakable,boxrule=0.6pt,arc=1.5mm,left=2mm,right=2mm,top=0.8mm,bottom=0.8mm,
        fonttitle=\bfseries\small,before skip=0.6em,after skip=0.6em}
\newtcolorbox{adire}[1][]{colback=violet!4,colframe=violet!60!black,
  title={À dire\ifx&#1&\else~— #1\fi}}
\newtcolorbox{attention}[1][]{colback=orange!4,colframe=orange!70!black,
  title={Attention\ifx&#1&\else~— #1\fi}}
\newtcolorbox{preuve}[1][]{colback=green!3,colframe=green!35!black,
  title={Nos chiffres\ifx&#1&\else~— #1\fi}}

\newcommand{\vz}{\bm{z}}
\newcommand{\va}{\bm{a}}
\newcommand{\vh}{\bm{h}}
\newcommand{\py}[1]{\texttt{#1}}
\newcommand{\F}{\mathsf{F}}
\newcommand{\tp}{^{\mathsf{T}}}

\begin{document}

{\centering
{\Large\bfseries\color{blue!40!black} Préparer la discussion avec M. Le Maître\par}
\vspace{0.15em}
{\large Plan de réponse, argumentaire, et lecture de sa note du 29/10\par}
\vspace{0.2em}
{\small Note interne (Adrian, Antoine) · 30/10/2026\par}}

\vspace{0.4em}
\textbf{En trois lignes.}
\begin{enumerate}
  \item \textbf{Sa note décrit notre architecture} : un GP des efforts, une VPP nourrie par sa moyenne, des
        sensibilités de la vitesse aux efforts, des calculs choisis pour la décision de forme. Nous parlons
        de la même chose ; c'est notre mail qui était flou.
  \item \textbf{Il va plus loin sur le critère} : il choisit les calculs pour réduire le \emph{regret} de
        conception (knowledge gradient sur les \emph{différences} de vitesse entre formes). Nous réduisons
        la variance des vitesses elles-mêmes. C'est la piste la plus intéressante de la réunion.
  \item \textbf{Sur les hyperparamètres}, il propose la validation croisée LOO : nous pouvons dire oui avec
        des arguments précis (§\,4).
\end{enumerate}

% ---------------------------------------------------------------------------------------------
\section{Sa note en une page}
% ---------------------------------------------------------------------------------------------
\begin{itemize}
  \item \textbf{§\,1–3 : le problème.} Trois objets distincts : la forme $u_f$ (conception), les réglages
        $u_r$ (adaptés en navigation) et l'état $a$ (fixé par l'équilibre). Le VPP résout
        $V^\star(u_f,v_w;\F)=\max_y V(a)$ sous $R=0$, $g\le0$, avec $y=(u_r,a)$. Le GP prédit le torseur
        $\F(u_f,z)$ en fonction de variables réduites $z=\Phi(u_r,a)=(V_f,h_f,\theta_1,\theta_2,\theta_3)$ :
        vitesse, enfoncement, trois angles.
  \item \textbf{§\,4–6 : le modèle et la sensibilité.} VPP sur la moyenne $\mu_N$ (comme nous). La
        sensibilité de la vitesse \emph{optimisée} aux efforts est donnée par les multiplicateurs de
        Lagrange du VPP : $b_N=B\tp\bar\lambda_N$, $\delta J\simeq b_N\tp\,\delta\F(\bar x_N)$, éq.~(5.14).
        Les réglages sont déjà réoptimisés dans cette formule (théorème de l'enveloppe).
  \item \textbf{§\,7 : la loi jointe des formes candidates.} $\tilde J_i=m_i+b_i\tp[\F(x^\star_i)-\mu_N(x^\star_i)]$,
        gaussienne, de covariance $S_{ij}=b_i\tp K_N(x^\star_i,x^\star_j)\,b_j$, éq.~(7.4). \textbf{Ce qui
        compte pour choisir est la différence} $\tilde J_i-\tilde J_j$, de variance $S_{ii}+S_{jj}-2S_{ij}$.
  \item \textbf{§\,8–10 : le critère.} Le regret espéré $r_N=\mathbb E_N[\max_i \tilde J_i]-\max_i m_i$, en
        nœuds. Le calcul suivant maximise sa baisse attendue (knowledge gradient, Frazier et al.~\cite{FPD09}),
        éq.~(9.11), calculable par Monte-Carlo en dimension 6. Pour deux formes : $\tilde\alpha=H(|d|,\tau(x))$,
        éq.~(10.6). Une incertitude commune à toutes les formes ne vaut rien.
  \item \textbf{§\,11 : la boucle}, les contrôles et un \textbf{critère d'arrêt en nœuds} : $\tilde r_N\le
        \varepsilon_V$, avec des bornes simples (11.2).
  \item \textbf{§\,12 : plusieurs conditions}, moyenne pondérée $\sum_k\omega_k V^\star_k$ (notre profil
        d'usage), avec \emph{toutes} les covariances croisées entre conditions, éq.~(12.6)–(12.7).
\end{itemize}

% ---------------------------------------------------------------------------------------------
\section{Correspondance entre sa note et notre code}
% ---------------------------------------------------------------------------------------------
{\small
\begin{tabularx}{\linewidth}{@{}>{\raggedright\arraybackslash}p{3.0cm}>{\raggedright\arraybackslash}X>{\raggedright\arraybackslash}p{4.4cm}@{}}
\toprule
Sa note & Chez nous & Dans le code \\
\midrule
$u_f$ (forme) & $\vz\in\mathbb R^{10}$, scores PCA des 16 poids de Kulfan & \py{SectionPCA} \\
$u_r$ (réglages) & safran, \enquote{flat} (dépuissance) ; le cant du foil est \emph{fixé} (à confirmer avec Antoine) & VPP MDS, \py{SOLVE\_VARS} \\
$a$ (état) & vitesse, dérive, gîte, assiette, déplacement de coque & idem \\
$z=\Phi(u_r,a)$ & $\va=$ (cant, rake, yaw) \emph{effectifs} : 3 angles seulement & \py{FOIL\_INPUTS} \\
$\F$ ($q=6$) & 7 co-krigeages : torseur + $C_{p,\min}$, en \emph{coefficients} & \py{TorsorCoKriging} \\
$B$ (repères, bras de levier) & fait par MDS : notre modèle remplace la surface de réponse du foil & \py{SurrogateFoilRS} \\
$\mu_N$, $\bar J_N$ & VPP avec la moyenne des co-krigeages & \py{IntensSYVPP} \\
$b_N=B\tp\bar\lambda_N$ & $d_k=\partial V^\ast/\partial C_k$ \textbf{par différences finies}, VPP complète réoptimisée & sensibilités du titulaire \\
$\omega_k$, $v_w^{(k)}$ & $w_c$, nœuds du profil d'usage & \py{OperatingProfile} \\
$\mathcal D_f^c$ (candidats) & titulaire + 3 concurrentes du criblage, optima des régions & \py{VPPSectionOptimizer} \\
\bottomrule
\end{tabularx}}

\textbf{Ce qui coïncide} (à dire : nous avons construit, empiriquement, une version de son cadre) :
\begin{itemize}
  \item un GP des efforts partagé par toutes les formes et toutes les conditions (son § 12 : \enquote{un seul GP,
        pas un GP par vent}) ;
  \item la VPP sur la moyenne du GP, avec la mise en garde $\mathbb E_N[J]\ne J(\mu_N)$ (sa (5.3)) ;
  \item \textbf{nos $d_k$ sont ses $b_N$} : décaler une sortie du modèle et refaire la VPP réoptimisée, c'est
        exactement le contrôle par différences finies de sa §\,11.4. Il le propose comme vérification ; c'est
        notre méthode de calcul ;
  \item les calculs choisis pour la vitesse, pondérés par la sensibilité, et non pour la précision des efforts
        (sa §\,5.2 : \enquote{une somme non pondérée des variances des six efforts n'est pas un indicateur de
        précision sur la performance}) ;
  \item le coût dans le choix : son score $\tilde\alpha/c(x)$ (§\,9.5) est notre rendement gain/coût ;
  \item le $C_{p,\min}$ comme sortie à part : il note qu'une contrainte de cavitation ne se déduit pas du
        torseur (§\,3.1) ;
  \item les VPP à départs multiples, contre les branches locales (sa §\,6.3).
\end{itemize}

\textbf{Ce qui diffère, et ce que cela nous apprend :}
\begin{enumerate}
  \item \textbf{Le critère} (le point principal). Notre ALC réduit $\sum_c w_c\,\mathrm{Var}\,V_c$ autour des
        états : une somme de variances de vitesses, forme par forme. Il ne regarde ni les \emph{différences}
        entre formes, ni l'écart de leurs moyennes. Son §\,9.2 le dit : \enquote{une somme arbitraire de
        réductions de variance n'est pas identique à la réduction du regret}. Concrètement, notre critère
        dépense des calculs pour une incertitude commune au titulaire et à ses concurrentes, qui ne change
        aucun classement.
  \item \textbf{Les covariances croisées.} Au criblage, notre $\sigma$ ne garde que les variances (pas
        $-2S_{ij}$) et suppose les conditions indépendantes. Sa (7.6) et sa (12.6) gardent tout. Nos
        concurrentes sont proches du titulaire, donc très corrélées : notre $\sigma$ de \emph{différence} est
        probablement trop grand.
  \item \textbf{Les sensibilités par multiplicateurs.} $b_N=B\tp\bar\lambda$ sort du solveur sans VPP
        supplémentaire, alors que nos différences finies coûtent 12 solves VPP par condition. Notre SLSQP
        (scipy 1.14 dans \py{.venv-mds}) ne renvoie pas, à notre connaissance, les multiplicateurs : à vérifier.
  \item \textbf{Les variables réduites.} Il a 5 variables ($V_f$, $h_f$ et 3 angles), nous 3 angles. Nos
        efforts sont des coefficients (la vitesse n'entre que par la pression dynamique et la correction
        FSI), et la hauteur de vol (\py{Hvol}) a été retirée des entrées des vrais codes, faute de surface
        libre dans AVL et XFOIL. Ce sont ses hypothèses H3 et H5, que nous devons écrire explicitement.
  \item \textbf{L'arrêt.} Nous arrêtons quand 10 mises à jour gagnent moins de 0,01 nd : c'est une
        heuristique. Son $\tilde r_N\le\varepsilon_V$ est en nœuds, avec des bornes calculables (11.2).
  \item \textbf{Le multi-fidélité} n'est pas dans sa note (un seul niveau). Le co-krigeage donne les mêmes
        objets ($\mu_N$, $K_N$ du niveau haut) ; le coût par niveau entre par son $\tilde\alpha/c$.
  \item \textbf{Les lots.} Son critère est à un pas ; nous choisissons 16 calculs à la fois (kriging believer).
        Il faudrait une version par lot (q-KG).
\end{enumerate}

\subsection*{Quatre différences, dans ses notations et nos conventions}

\textbf{a) Une sensibilité par forme candidate.} Chez lui, chaque forme candidate $u_f^{(i)}$ a son propre VPP
nominal, son point d'équilibre $x^\star_i=(u_f^{(i)},\Phi(\bar y_i))$ et sa propre sensibilité
$b_i=B(\bar y_i)\tp\bar\lambda_i$ (7.2). Sa vitesse linéarisée est $\tilde J_i=m_i+b_i\tp[\F(x^\star_i)-\mu_N(x^\star_i)]$.
Chez nous, $d_{kc}$ n'est calculé qu'à la titulaire $\vz_0$, à ses états, puis servi à tout le monde :
\begin{itemize}
  \item au criblage, une section $\vz$ est jugée avec les $d_{kc}$ de $\vz_0$ \emph{et à l'attitude $\va_c$ de
        $\vz_0$}, pas à son propre équilibre ;
  \item dans l'ALC, les points autour des états des 3 concurrentes pèsent $w_c\,d_{kc}^2$ de la titulaire.
\end{itemize}
C'est juste tant que les concurrentes sont proches de la titulaire. C'est faux si leur équilibre diffère : une
concurrente en butée à 24 nd a une sensibilité nulle à cette condition, la titulaire non.

\textbf{b) Les sensibilités par les multiplicateurs.} Son VPP (3.4) est $\max_y V(a)$ sous $R(y)=0$ et
$g(y)\le0$. Chez nous, $y$ sont les 7 inconnues de SLSQP, $R$ les 6 équations d'équilibre et $g$ leurs bornes.
\begin{itemize}
  \item À l'optimum, chaque équation $j$ a un multiplicateur $\bar\lambda_j$ : la vitesse gagnée par unité
        d'effort ajoutée dans cette équation (un \enquote{prix ombre}, en nd/N ou nd/(N·m)).
  \item $B_{jk}=\partial R_j/\partial C_k$, à $y$ fixé, dit comment le coefficient $k$ du foil entre dans
        l'équation $j$ : pression dynamique, surface, repère, bras de levier.
  \item Alors $d_k=(B\tp\bar\lambda)_k=\sum_j\bar\lambda_j B_{jk}$. Un seul solve donne les 6 $d_k$, au lieu de
        nos 12 solves par condition (pour chaque $k$, la base refaite dans le même processus et le solve
        décalé de 3 \%).
  \item \textbf{En pratique} : si SLSQP ne renvoie pas $\bar\lambda$, on le reconstruit à l'optimum par la
        stationnarité $\nabla_yV+A\tp\bar\lambda=0$, sur les inconnues qui ne sont pas en butée, avec
        $A=\partial R/\partial y$ par différences finies sur le résidu (des évaluations, pas des solves). $B$
        s'obtient de même par 6 évaluations du résidu. Attention aux conventions (sa 5.17) : notre SLSQP minimise
        $-V$ sous des résidus divisés par \py{EQ\_SCALE}.
\end{itemize}

\textbf{c) L'arrêt en nœuds.} Chez lui, avec les candidates (titulaire, concurrentes, optima) et leur loi
jointe $\tilde{\bm J}\sim\mathcal N(\bm m,S)$, le regret espéré $\tilde r_N=\mathbb E[\max_i\tilde J_i]-\max_i m_i$
est la vitesse qu'on s'attend à perdre en choisissant aujourd'hui. On s'arrête quand $\tilde r_N\le\varepsilon_V$
(par exemple 0,005 nd), avec les bornes (11.2). Chez nous, une région s'arrête quand ses 10 dernières mises à
jour ont gagné moins de 0,01 nd. \textbf{Le nôtre dit que la recherche ne progresse plus ; le sien, que le choix
est sûr.} Une recherche peut stagner avec un choix encore incertain, ou progresser alors que tout est joué.

\textbf{d) Les entrées du GP.} Chez lui, $\F(u_f,z)$ avec $z=\Phi(u_r,a)=(V_f,h_f,\theta_1,\theta_2,\theta_3)$ ; ses
hypothèses H3 et H5 disent que ces 5 variables suffisent. Chez nous, $\vh(\vz,\va)$, où $\va$ sont ses
$\theta_1,\theta_2,\theta_3$. Il nous manque $V_f$ et $h_f$ :
\begin{itemize}
  \item nos sorties sont des coefficients, l'effort vaut $\tfrac12\rho V_f^2 S\,C_k$ : on suppose les
        coefficients indépendants de la vitesse (pas d'effet de Reynolds, bdfoil\_v1 calcule à vitesse fixe),
        sauf la déformation (la correction FSI prend la vitesse) ;
  \item aucun effet de surface libre : AVL et XFOIL n'en ont pas, et \py{Hvol} (hauteur de vol) a été retiré
        des entrées des vrais codes.
\end{itemize}
Ce sont des hypothèses à écrire, pas des erreurs. Elles deviendront fausses si un niveau CFD avec surface libre
entre dans le modèle.

\begin{adire}[les changements à faire, dans cet ordre]
\begin{enumerate}
  \item \textbf{Les sensibilités par les multiplicateurs}, $d=B\tp\bar\lambda$, vérifiées contre nos différences
        finies (plusieurs pas, différence centrée). Elles rendent abordable une sensibilité $b_i$ par candidate.
  \item \textbf{La loi jointe complète}, $S_{ij}$ entre formes et entre conditions (sa 12.6), et le regret
        $\tilde r_N$, calculés et écrits à chaque mise à jour : un diagnostic, sans changer la boucle.
  \item \textbf{L'arrêt sur $\tilde r_N\le\varepsilon_V$} au lieu de 0,01 nd en 10 mises à jour. On garde les
        redémarrages : le regret ne voit que les candidates connues.
  \item \textbf{Le critère de regret (9.11) à la place de l'ALC}, en version par lot et à deux niveaux, comparé
        à l'ALC sur un cas test avant les vrais runs.
  \item \textbf{La calibration de $\sigma$} par LOO (§\,4) : le regret dépend directement de $S$.
  \item \textbf{Les hypothèses sur $V_f$ et $h_f$}, écrites ; $h_f$ en entrée si un niveau avec surface libre
        arrive.
\end{enumerate}
\end{adire}

% ---------------------------------------------------------------------------------------------
\section{Les points 2 et 3 de sa réponse : reformuler}
% ---------------------------------------------------------------------------------------------
Il évoquait la thèse de Malo (des GP \enquote{en interaction via une équation d'équilibre}). Sa note montre que
le cadre qu'il a en tête est bien le nôtre. Le message à redire, dans ses notations :
\begin{itemize}
  \item \textbf{le GP apprend $\F(u_f,z)$ à partir de calculs à entrées imposées} (sa (4.1)) ; aucun GP n'est
        appris à travers l'équilibre ;
  \item \textbf{l'équilibre n'intervient qu'en aval} : le VPP nominal sur $\mu_N$, la sensibilité $b_N$, et le
        choix des calculs ;
  \item \textbf{notre mail était inexact} sur trois points : $f(x)=g(h(x))$ au lieu de $\F(u_f,\Phi(y))$ avec
        $y$ choisi par le VPP ; 6 co-krigeages au lieu de 7 ; des \enquote{angles d'orientation
        géométrique} qui sont en fait ses $\theta_j$, des variables de fonctionnement.
\end{itemize}

% ---------------------------------------------------------------------------------------------
\section{Le point 1 : MAP, MLE et validation croisée}
% ---------------------------------------------------------------------------------------------
\begin{adire}[notre position]
Oui à la LOO : d'abord comme \textbf{diagnostic} sur nos modèles, puis comme \textbf{critère d'estimation}
des hyperparamètres, comparée au MAP sur les mêmes cas (Hartmann 6D, section 2D, données des runs).
\end{adire}

\begin{enumerate}
  \item \textbf{Sur le fond, il a raison} : le prior est une régularisation, assumée comme un choix pragmatique.
        Nos priors sont faibles, sur des entrées normalisées dans $[0,1]$ (Inverse-Gamma(3, 2) sur les
        longueurs, de mode 0,5). Ils servent surtout à éviter les longueurs collées aux bornes.
  \item \textbf{La LOO coûte peu} : formules fermées de Dubrule~\cite{Dub83},
        $y_i-\hat y_{-i}=[K^{-1}\bm y]_i/[K^{-1}]_{ii}$ et $\sigma^2_{-i}=1/[K^{-1}]_{ii}$, une seule
        factorisation.
  \item \textbf{Les gradients existent} (il ne les avait pas sous la main) : Rasmussen \& Williams~\cite{RW06}
        (§\,5.4.2), d'après Sundararajan \& Keerthi~\cite{SK01}.
  \item \textbf{Pour le co-krigeage, c'est dans Le Gratiet \& Cannamela~\cite{LGC15}}, l'article dont part notre
        critère (\enquote{using fast cross-validation techniques}).
  \item \textbf{La nuance qui joue pour nous} : la LOO en erreur quadratique ne fixe pas $\sigma^2$. Or son
        critère comme le nôtre est fait de covariances ($S$, $C_N$, $\Lambda_N$) : leur calibration compte
        autant que la moyenne. Il faut la log-densité prédictive LOO, ou $\sigma^2$ calée sur les erreurs LOO
        réduites (Bachoc~\cite{Bac13}, qui montre aussi la robustesse de la validation croisée quand la
        covariance est mal spécifiée).
\end{enumerate}

\begin{preuve}[ce que dit la LOO du modèle du 8\textsuperscript{e} run (30/10)]
\textbf{La LOO en forme fermée} de mfego (tous les niveaux d'un point retirés ensemble), hyperparamètres du run,
899 points BF et 276 points HF :
\begin{itemize}
  \item \textbf{BF, les 7 sorties} : $Q^2\ge0{,}97$, écart type des erreurs réduites $z_{\text{rms}}$ de 0,97 à
        1,07, couverture à 95 \% de 0,93 à 0,95 : \textbf{bien calibré} ;
  \item \textbf{HF, torseur} : $Q^2\ge0{,}999$, $z_{\text{rms}}$ de 0,77 à 0,92, couverture 0,96 à 0,99 : la
        variance est un peu \textbf{trop grande} (environ 25 \% sur l'écart type), donc prudente ;
  \item \textbf{HF, $C_{p,\min}$} : $Q^2=0{,}56$, $z_{\text{rms}}=1{,}18$, couverture 0,91 : mal prédit et trop
        sûr de lui. Sans effet sur la vitesse, mais à revoir pour une contrainte de cavitation ;
  \item \textbf{en vitesse} (erreurs LOO du torseur × sensibilités de l'état final le plus proche) : erreur RMS
        0,015 nd sur les 276 points HF, $|z|\le2$ sur 97 \% ; 0,001 nd près de la meilleure section.
\end{itemize}
\textbf{Mais la validation de fin de run dit autre chose loin des données} : aux meilleures sections des runs 5
et 6, l'erreur contre AVL × XFOIL est de 0,013 nd et $|z|\le2$ sur 65 à 69 \% des états seulement, soit une
variance sous-estimée d'un facteur 2 environ.

\textbf{La lecture} : la LOO teste l'\emph{interpolation} (le point retiré a des voisins dans son nuage) ; elle ne
voit pas l'\emph{extrapolation}, là où la variance est fausse. Régler les hyperparamètres par LOO ne corrigerait
donc pas ce défaut. C'est un point à lui soumettre : avec son critère de regret, une variance sous-estimée loin
des données fait croire un classement plus sûr qu'il n'est.
\end{preuve}

\begin{attention}[ne pas survendre le \enquote{peu de données}]
Le \enquote{$n<20$} de notre mail vaut pour les études de test et le début d'un run. À la fin du
8\textsuperscript{e} run, il y a 899 points BF et 276 points HF en dimension 13 : le prior pèse alors peu, sauf
sur $\rho$.
\end{attention}

% ---------------------------------------------------------------------------------------------
\section{Les questions qu'il risque de poser}
% ---------------------------------------------------------------------------------------------
{\small
\begin{tabularx}{\linewidth}{@{}>{\raggedright\arraybackslash}p{5.0cm}>{\raggedright\arraybackslash}X@{}}
\toprule
Question & Réponse honnête \\
\midrule
Où est l'enfoncement $h_f$ ? Et $V_f$ ? & Pas d'entrée : coefficients (vitesse par la pression dynamique et la FSI), pas de surface libre dans AVL/XFOIL, \py{Hvol} retiré pour les vrais codes. Ce sont des hypothèses à écrire. \\
Vos sensibilités tiennent-elles compte des réglages ? & Oui : la VPP est refaite \emph{optimisée} après le décalage. C'est sa (5.14) par différences finies, son test de la §\,11.4. \\
Mettre la moyenne dans la VPP est biaisé. & Oui, au second ordre ; contrôlé en fin de run (0,0006 nd à l'optimum du 8\textsuperscript{e} run). \\
Pourquoi 7 GP indépendants ? & Simplicité. Il le dit lui-même : l'indépendance entre composantes n'autorise pas à oublier les covariances entre entrées (sa §\,4.1), que l'ALC garde. \\
Et les autres branches d'équilibre ? & VPP à départs multiples ; une condition sans équilibre ne pèse rien. Sa (6.4) est plus propre. \\
L'optimum est-il global ? & Non garanti : redémarrages. Au 8\textsuperscript{e} run, deux régions à 13,566 et 13,560 nd : un plateau. Avec sa (7.3), c'est un regret petit, donc un bon signe. \\
\bottomrule
\end{tabularx}}

% ---------------------------------------------------------------------------------------------
\section{Ce que nous pouvons lui demander}
% ---------------------------------------------------------------------------------------------
\begin{enumerate}
  \item Le knowledge gradient avec \textbf{deux niveaux de fidélité} : faut-il conditionner sur la HF seule,
        ou sur le couple (BF, HF) comme notre ALC ?
  \item Une version \textbf{par lot} de (9.11), puisque nos calculs tournent par 16.
  \item Comment \textbf{renouveler l'ensemble de candidats} $\mathcal D_f^c$ (sa §\,11.2) : notre criblage
        optimiste lui convient-il ?
  \item Son expérience de la LOO : erreur quadratique ou log-densité ? Comment caler $\sigma^2$ ?
  \item Ce qui bloquait dans la thèse de Malo, pour savoir ce qu'il faut éviter.
\end{enumerate}

% ---------------------------------------------------------------------------------------------
\section{Le déroulé proposé pour la réunion (45 min)}
% ---------------------------------------------------------------------------------------------
\begin{enumerate}
  \item \textbf{5 min :} la correspondance de la §\,2, dans ses notations : \enquote{nous avons construit une
        version empirique de votre cadre}.
  \item \textbf{10 min :} notre boucle (\py{structure\_vpp\_acceleration.pdf}, figure~1), en marquant ce qui
        coïncide et ce qui diffère.
  \item \textbf{10 min :} les résultats du 8\textsuperscript{e} run, avec ses figures (\py{convergence},
        \py{sections}, \py{validation}, \py{profile\_profile\_v2}). +0,256 nd sur la section actuelle
        (+1,9 \%), en 2 h 49, 1\,764 calculs de codes (1\,175 gardés dans le plan final).
  \item \textbf{15 min :} le critère de regret et l'ordre proposé pour l'adopter (§\,2), avec ses conseils.
  \item \textbf{5 min :} la LOO (§\,4).
\end{enumerate}

\begin{thebibliography}{9}\footnotesize\setlength{\itemsep}{0pt}
\bibitem{FPD09} P. Frazier, W. Powell, S. Dayanik, \emph{The knowledge-gradient policy for correlated normal
beliefs}, INFORMS Journal on Computing 21(4), 599–613, 2009 (sa référence [3]).
\bibitem{Dub83} O. Dubrule, \emph{Cross validation of kriging in a unique neighborhood}, Mathematical Geology
15(6), 687–699, 1983.
\bibitem{RW06} C. E. Rasmussen, C. K. I. Williams, \emph{Gaussian Processes for Machine Learning}, MIT Press,
2006.
\bibitem{SK01} S. Sundararajan, S. S. Keerthi, \emph{Predictive approaches for choosing hyperparameters in
Gaussian processes}, Neural Computation 13(5), 1103–1118, 2001.
\bibitem{Bac13} F. Bachoc, \emph{Cross validation and maximum likelihood estimations of hyper-parameters of
Gaussian processes with model misspecification}, Computational Statistics \& Data Analysis 66, 55–69, 2013.
\bibitem{LGC15} L. Le Gratiet, C. Cannamela, \emph{Cokriging-based sequential design strategies using fast
cross-validation techniques for multi-fidelity computer codes}, Technometrics 57(3), 418–427, 2015.
\end{thebibliography}

\end{document}
